% Mathématiques 6e — cours V3
% Source LaTeX rendue côté Astro/KaTeX.

\section{Objectifs}

Dans ce chapitre, on apprend à :

\begin{itemize}
\item distinguer un point, un segment, une droite et une distance ;
\item utiliser la notation d’une distance entre deux points ;
\item caractériser et construire le milieu d’un segment ;
\item distinguer cercle et disque ;
\item maîtriser les mots rayon, diamètre et corde ;
\item traduire une position par une condition sur une distance au centre ;
\item construire une médiatrice ;
\item utiliser la propriété d’équidistance de la médiatrice ;
\item résoudre des problèmes de localisation géométrique.
\end{itemize}

\section{1. Distance entre deux points}

Soient deux points $A$ et $B$.

Le segment ayant pour extrémités $A$ et $B$ se note :

\[
[AB].
\]

La \textbf{distance} entre $A$ et $B$ est la longueur de ce segment.

Elle se note :

\[
AB.
\]

Ainsi, si le segment $[AB]$ mesure :

\[
4{,}5\ \text{cm},
\]

on écrit :

\[
AB=4{,}5\ \text{cm}.
\]

La notation $[AB]$ désigne un objet géométrique ; la notation $AB$ désigne une longueur.

\section{2. Une distance est toujours positive ou nulle}

Pour deux points quelconques :

\[
AB\ge0.
\]

Et :

\[
AB=0
\]

uniquement lorsque les deux points sont confondus.

De plus :

\[
AB=BA.
\]

L’ordre dans lequel on nomme les deux points ne change pas la distance.

\section{3. Le milieu d’un segment}

Un point $I$ est le milieu du segment $[AB]$ lorsqu’il vérifie deux conditions :

\begin{enumerate}
\item $I$ appartient au segment $[AB]$ ;
\item les deux longueurs sont égales :
\end{enumerate}

\[
AI=IB.
\]

Si :

\[
AB=8\ \text{cm},
\]

alors le milieu $I$ vérifie :

\[
AI=IB=4\ \text{cm}.
\]

La seule égalité des distances ne suffit pas à définir le milieu : il faut également que $I$ soit situé sur le segment.

\section{4. Exemple développé 1 — retrouver une longueur}

On sait que $I$ est le milieu de $[AB]$ et que :

\[
AI=3{,}7\ \text{cm}.
\]

Comme $I$ est le milieu :

\[
IB=AI=3{,}7\ \text{cm}.
\]

Donc :

\[
AB=AI+IB.
\]

Ainsi :

\[
AB=3{,}7+3{,}7=7{,}4\ \text{cm}.
\]

On obtient :

\[
\boxed{AB=7{,}4\ \text{cm}}.
\]

\section{5. Le cercle}

Un cercle est défini par :

\begin{itemize}
\item un centre ;
\item un rayon.
\end{itemize}

Le cercle de centre $O$ et de rayon $r$ est l’ensemble des points $M$ qui vérifient :

\[
\boxed{OM=r}.
\]

Tous les points du cercle sont donc situés à la même distance du centre.

\section{6. Le disque}

Le disque de centre $O$ et de rayon $r$ contient tous les points $M$ tels que :

\[
\boxed{OM\le r}.
\]

Il comprend :

\begin{itemize}
\item les points strictement à l’intérieur :
\end{itemize}

\[
OM<r;
\]

\begin{itemize}
\item les points du cercle frontière :
\end{itemize}

\[
OM=r.
\]

Le cercle est donc le contour du disque.

\section{7. Vocabulaire du cercle}

\coursefigure{/images/mathematiques/6eme/geometrie-cercle-disque.svg}{Disque avec centre, rayon, diamètre et corde représentés.}{Un rayon relie le centre au cercle ; un diamètre est une corde qui passe par le centre.}

Si $O$ est le centre et $A$ un point du cercle, le segment $[OA]$ est un rayon.

Sa longueur est :

\[
OA=r.
\]

Un diamètre relie deux points du cercle en passant par le centre.

Si $[AB]$ est un diamètre :

\[
AB=2r.
\]

Une corde relie deux points du cercle sans être nécessairement un diamètre.

\section{8. Rayon et diamètre}

La relation entre rayon et diamètre est :

\[
\boxed{d=2r}.
\]

Et :

\[
\boxed{r=\dfrac{d}{2}}.
\]

Par exemple, pour :

\[
r=3{,}5\ \text{cm},
\]

le diamètre vaut :

\[
d=2\times3{,}5=7\ \text{cm}.
\]

\section{9. Position d’un point par rapport à un disque}

Pour un disque de centre $O$ et de rayon $r$, on compare $OM$ à $r$.

\subsection{Point à l’intérieur}

\[
OM<r.
\]

\subsection{Point sur le cercle}

\[
OM=r.
\]

\subsection{Point à l’extérieur}

\[
OM>r.
\]

Une simple comparaison de distances permet donc de déterminer la position du point.

\section{10. Exemple développé 2 — classer plusieurs points}

Un disque a pour centre $O$ et pour rayon :

\[
5\ \text{cm}.
\]

On connaît :

\[
OA=3\ \text{cm},
\]

\[
OB=5\ \text{cm},
\]

\[
OC=6{,}2\ \text{cm}.
\]

Comme :

\[
3<5,
\]

le point $A$ est à l’intérieur du disque.

Comme :

\[
5=5,
\]

le point $B$ appartient au cercle.

Comme :

\[
6{,}2>5,
\]

le point $C$ est à l’extérieur du disque.

\section{11. Cercle comme lieu géométrique}

Supposons que l’on cherche \textbf{tous les points situés exactement à une distance donnée d’un point $O$}.

Si la distance imposée est :

\[
4\ \text{cm},
\]

l’ensemble des solutions est le cercle de centre $O$ et de rayon :

\[
4\ \text{cm}.
\]

On dit que le cercle constitue un \textbf{lieu géométrique} : il représente l’ensemble de tous les points qui satisfont une condition.

\section{12. Disque comme zone autorisée}

Si l’on cherche tous les points situés à une distance \textbf{inférieure ou égale} à :

\[
4\ \text{cm}
\]

du point $O$, l’ensemble des solutions est le disque.

La différence entre :

\[
OM=4\ \text{cm}
\]

et :

\[
OM\le4\ \text{cm}
\]

est donc fondamentale.

\section{13. Médiatrice d’un segment}

La médiatrice du segment $[AB]$ est la droite qui :

\begin{itemize}
\item est perpendiculaire à la droite $(AB)$ ;
\item passe par le milieu de $[AB]$.
\end{itemize}

\coursefigure{/images/mathematiques/6eme/geometrie-mediatrice.svg}{Segment horizontal et droite médiatrice perpendiculaire passant par son milieu.}{La médiatrice de $[AB]$ est perpendiculaire à $(AB)$ et passe par le milieu du segment.}

Ces deux conditions appartiennent à la définition.

\section{14. Construction de la médiatrice au compas}

Pour construire la médiatrice de $[AB]$ :

\begin{enumerate}
\item ouvrir le compas avec un rayon suffisamment grand ;
\item tracer deux arcs de même rayon centrés en $A$ et en $B$ ;
\item obtenir deux points d’intersection $M$ et $N$ ;
\item tracer la droite $(MN)$.
\end{enumerate}

Comme les arcs ont le même rayon :

\[
MA=MB
\]

et :

\[
NA=NB.
\]

Les deux points $M$ et $N$ sont donc équidistants de $A$ et $B$.

\section{15. Propriété caractéristique de la médiatrice}

Un point $M$ appartient à la médiatrice de $[AB]$ \textbf{si et seulement si} :

\[
\boxed{MA=MB}.
\]

Cette phrase contient deux sens de raisonnement.

\subsection{Sens direct}

Si $M$ est sur la médiatrice de $[AB]$, alors :

\[
MA=MB.
\]

\subsection{Sens réciproque}

Si :

\[
MA=MB,
\]

alors $M$ appartient à la médiatrice de $[AB]$.

Cette propriété est très puissante pour résoudre des problèmes.

\section{16. Exemple développé 3 — utiliser l’équidistance}

On sait qu’un point $M$ appartient à la médiatrice de $[AB]$ et que :

\[
MA=6{,}4\ \text{cm}.
\]

La propriété donne immédiatement :

\[
MB=MA.
\]

Donc :

\[
\boxed{MB=6{,}4\ \text{cm}}.
\]

Aucune mesure supplémentaire n’est nécessaire.

\section{17. Reconnaître qu’un point appartient à une médiatrice}

Supposons :

\[
MA=5\ \text{cm}
\]

et :

\[
MB=5\ \text{cm}.
\]

Alors :

\[
MA=MB.
\]

La propriété réciproque permet de conclure que :

\[
M
\]

appartient à la médiatrice de :

\[
[AB].
\]

On peut donc reconnaître un lieu géométrique à partir de distances.

\section{18. Problème de deux villages}

Deux villages sont représentés par les points $A$ et $B$.

On cherche tous les endroits situés à la même distance des deux villages.

La condition est :

\[
MA=MB.
\]

L’ensemble de ces points est donc :

\[
\boxed{\text{la médiatrice de }[AB]}.
\]

Cette modélisation peut servir, par exemple, à chercher des emplacements équidistants de deux sites.

\section{19. Croiser plusieurs conditions}

Une position géométrique peut devoir satisfaire plusieurs contraintes.

On cherche un point $M$ tel que :

\[
MA=MB
\]

et :

\[
MC=4\ \text{cm}.
\]

La première condition impose que $M$ soit sur la médiatrice de $[AB]$.

La seconde impose que $M$ soit sur le cercle de centre $C$ et de rayon :

\[
4\ \text{cm}.
\]

Les solutions sont donc les points d’intersection entre ces deux lieux géométriques.

Il peut y avoir :

\begin{itemize}
\item deux solutions ;
\item une seule solution ;
\item aucune solution.
\end{itemize}

\section{20. Exemple développé 4 — intersection de lieux}

On cherche un point situé :

\begin{itemize}
\item à égale distance de $A$ et $B$ ;
\item exactement à :
\end{itemize}

\[
3\ \text{cm}
\]

de $C$.

On construit :

\begin{enumerate}
\item la médiatrice de $[AB]$ ;
\item le cercle de centre $C$ et de rayon :
\end{enumerate}

\[
3\ \text{cm}.
\]

Les points d’intersection vérifient simultanément les deux conditions.

Le dessin ne remplace pas le raisonnement : chaque construction correspond à une propriété géométrique précise.

\section{21. Méthode — traduire une condition de distance}

\subsection{« À exactement une distance $r$ de $O$ »}

On construit un cercle :

\[
OM=r.
\]

\subsection{« À au plus une distance $r$ de $O$ »}

On considère un disque :

\[
OM\le r.
\]

\subsection{« À égale distance de $A$ et $B$ »}

On utilise la médiatrice :

\[
MA=MB.
\]

Traduire correctement la phrase est souvent l’étape la plus importante.

\section{22. Méthode — résoudre un problème de construction}

\begin{enumerate}
\item repérer les conditions données ;
\item traduire chaque condition par un objet géométrique ;
\item construire les objets avec règle et compas ;
\item rechercher leurs intersections ;
\item vérifier que chaque point obtenu satisfait toutes les conditions.
\end{enumerate}

Une construction doit être accompagnée d’une justification.

\section{23. Erreurs fréquentes}

\subsection{Confondre segment et distance}

Le segment s’écrit :

\[
[AB].
\]

La distance s’écrit :

\[
AB.
\]

\subsection{Confondre cercle et disque}

Le cercle correspond à :

\[
OM=r.
\]

Le disque correspond à :

\[
OM\le r.
\]

\subsection{Dire qu’un diamètre est deux rayons sans préciser l’alignement}

Un diamètre passe par le centre et a ses deux extrémités sur le cercle.

\subsection{Oublier une condition du milieu}

Pour que $I$ soit le milieu de $[AB]$, il faut :

\[
I\in[AB]
\]

et :

\[
AI=IB.
\]

\subsection{Croire que seul le milieu appartient à la médiatrice}

La médiatrice contient une infinité de points.

Tous vérifient :

\[
MA=MB.
\]

\subsection{Utiliser un arc de rayon différent dans la construction}

Les arcs centrés en $A$ et $B$ doivent être tracés avec la même ouverture du compas pour obtenir des points équidistants.

\section{24. À retenir}

La distance entre $A$ et $B$ est la longueur du segment $[AB]$ et se note :

\[
AB.
\]

Le milieu $I$ de $[AB]$ vérifie :

\[
I\in[AB]
\]

et :

\[
AI=IB.
\]

Le cercle de centre $O$ et de rayon $r$ est caractérisé par :

\[
\boxed{OM=r}.
\]

Le disque est caractérisé par :

\[
\boxed{OM\le r}.
\]

La médiatrice de $[AB]$ est la droite perpendiculaire à $(AB)$ passant par son milieu.

Sa propriété essentielle est :

\[
\boxed{
M\in\text{médiatrice de }[AB]
\iff
MA=MB
}.
\]

Ces objets permettent de transformer des contraintes de distance en constructions géométriques précises.
